\section{Formulación del problema y planificación del estudio}

\subsection{Descripción del sistema bajo estudio}

El sistema analizado comprende el flujo de pacientes pediátricos desde su
ingreso a la guardia hasta la finalización de la atención médica primaria.

Los pacientes pueden ingresar de manera peatonal por la entrada principal o
mediante ambulancia. En el primer caso atraviesan el proceso de admisión en
secretaría y luego son evaluados en triage. Los pacientes críticos que
ingresan en ambulancia son derivados directamente al Shock Room, sin pasar
por los procesos administrativos.

Durante el triage, una enfermera clasifica a los pacientes según cuatro
niveles de prioridad. Posteriormente, los pacientes esperan atención médica
de acuerdo con esa prioridad y con la disponibilidad de los consultorios. El
modelo considera cuatro consultorios médicos y el Shock Room.

\subsection{Formulación del problema}

La problemática principal se encuentra asociada a la posibilidad de que la
clasificación realizada durante el triage no coincida con la gravedad real
del paciente.

Si bien el proceso se basa en criterios clínicos establecidos, en la práctica
intervienen factores humanos relacionados con la experiencia, el criterio y
las condiciones de trabajo del profesional responsable de la evaluación. Esto
puede originar errores de clasificación que modifican artificialmente la
prioridad con la que un paciente ingresa al sistema de atención.

Se distinguen dos tipos de error:

\begin{itemize}
  \item \textbf{Sobre-triage:} un paciente cuya condición real presenta
  gravedad menor recibe una clasificación de mayor prioridad. Como
  consecuencia, accede anticipadamente a recursos que deberían estar
  disponibles para pacientes con urgencia real mayor.
  \item \textbf{Sub-triage:} un paciente cuya condición real presenta
  gravedad mayor recibe una clasificación de menor prioridad. En este caso
  permanece en la sala de espera más tiempo del adecuado antes de acceder a
  la atención médica.
\end{itemize}

Ambos fenómenos alteran el funcionamiento general de la guardia. El
sobre-triage incrementa innecesariamente la utilización de los recursos
destinados a pacientes prioritarios; el sub-triage produce demoras
particularmente relevantes en pacientes que necesitan atención rápida.

El problema de decisión que guía el estudio puede expresarse de la siguiente
manera:

\begin{quote}
\textit{¿En qué medida los errores de clasificación producidos durante el
triage afectan los tiempos de espera y permanencia de los pacientes en la
guardia, especialmente aquellos que presentan una urgencia real, y qué efecto
tendría sobre el desempeño general del sistema una modificación en la
capacidad del puesto de clasificación?}
\end{quote}

La simulación permite estudiar este problema sin intervenir sobre el
funcionamiento real de la guardia y observar el comportamiento del sistema
bajo distintas condiciones de demanda y clasificación.

\subsection{Objetivos del estudio}

\subsubsection*{Objetivo general}

Desarrollar un modelo de simulación de la guardia pediátrica del Hospital
Vilela mediante AnyLogic que permita analizar el impacto de los errores de
clasificación del triage sobre los tiempos de espera, la permanencia de los
pacientes y la utilización de los recursos de atención.

\subsubsection*{Objetivos específicos}

\begin{itemize}
  \item Representar el flujo principal de pacientes desde su ingreso hasta la
  finalización de la atención médica primaria.
  \item Incorporar la variabilidad existente en los arribos y en los tiempos
  de atención.
  \item Analizar los tiempos de espera de los pacientes según su nivel de
  urgencia.
  \item Evaluar las demoras experimentadas por pacientes cuya condición real
  requiere atención prioritaria.
  \item Analizar la congestión generada en la sala de espera y la utilización
  de los principales recursos del sistema.
  \item Comparar el funcionamiento actual con un escenario alternativo
  orientado a mejorar el desempeño.
  \item Determinar, mediante corridas múltiples y análisis estadístico, si las
  diferencias observadas entre escenarios resultan significativas.
\end{itemize}

\subsection{Alcance del modelo}

El modelo comienza con el arribo del paciente a la guardia --- por el ingreso
peatonal o por ambulancia --- y finaliza cuando, luego de recibir la atención
médica primaria, se determina uno de dos destinos finales: alta a domicilio
o internación.

La internación se considera únicamente como destino final dentro del modelo.
No se representa la permanencia posterior en las salas de internación ni los
procesos que ocurren dentro de ellas.

Tampoco se modelan de manera independiente las derivaciones internas hacia
traumatología, radiografías o medicación e inyecciones. Si bien estos
sectores forman parte de la guardia real, se excluyen para mantener un nivel
de abstracción adecuado al problema principal estudiado.

El flujo principal considerado es:

\begin{center}
Ingreso $\rightarrow$ Admisión $\rightarrow$ Triage $\rightarrow$ Espera según
prioridad $\rightarrow$ Atención médica $\rightarrow$ Alta / Internación
\end{center}

Para los pacientes críticos que ingresan mediante ambulancia se contempla un
flujo diferenciado que evita los procesos administrativos y de triage y
permite el acceso directo al Shock Room.

\begin{pendiente}
La exclusión de traumatología, radiografías y medicación retira capacidad de
atención del modelo pero no retira pacientes: la totalidad de los arribos se
dirige a los cuatro consultorios generales. Las ejecuciones piloto
(sección~\ref{sec:piloto}) muestran que esta decisión es la principal
candidata a explicar la saturación observada. El grupo debe resolver entre
dos alternativas: (a) declarar que la tasa de arribos representa únicamente
a los pacientes destinados a consultorio general, o (b) incorporar una
derivación explícita hacia esos sectores. Este informe documenta el modelo
tal como está implementado hoy, es decir, sin derivación.
\end{pendiente}

\subsection{Interacciones y disciplina de las colas}

Los pacientes que ingresan por la entrada principal son atendidos inicialmente
en secretaría mediante una fila común con disciplina FIFO. Posteriormente
pasan al proceso de triage, cuya cola previa también es FIFO.

Luego del triage la disciplina cambia. La selección de pacientes para los
consultorios se realiza de acuerdo con la prioridad clínica asignada:

\begin{center}
Rojo $\rightarrow$ Naranja $\rightarrow$ Amarillo $\rightarrow$ Verde
\end{center}

Por este motivo, un paciente que llega después puede ser atendido antes que
otro que se encuentra esperando, si posee un nivel de prioridad superior.

Corresponde señalar una característica relevante del sistema: \textbf{la cola
previa al triage es una cola ciega}. Antes de la clasificación nadie conoce la
gravedad del paciente, de modo que todo el tiempo de espera anterior al
triage transcurre necesariamente sin priorización alguna.

\subsection{Supuestos iniciales}

\begin{itemize}
  \item El análisis se concentra en un período de alta demanda de la guardia,
  correspondiente a un día pico invernal.
  \item Los arribos de pacientes son estocásticos y siguen un proceso de
  Poisson.
  \item Los tiempos asociados a los distintos servicios son variables
  aleatorias con distribución triangular.
  \item Los pacientes se clasifican en cuatro niveles de prioridad. La
  clasificación asignada durante el triage puede diferir de la gravedad real
  del paciente, por tratarse de una evaluación con componente subjetivo.
  \item Los pacientes ingresados mediante ambulancia con condición crítica
  acceden directamente a la atención de emergencia.
  \item Secretaría funciona mediante una cola FIFO; la atención médica
  posterior al triage utiliza una cola con prioridad según el código asignado.
  \item Los únicos destinos finales son la internación y el alta a domicilio.
  \item No se representan como procesos independientes las derivaciones a
  traumatología, radiografías o medicación.
  \item La probabilidad de error de clasificación es independiente de la carga
  de trabajo de la enfermera.
\end{itemize}

El último supuesto merece una aclaración. Al mantener constante la tasa de
error, el escenario alternativo modifica únicamente los tiempos de espera y
no la calidad de la clasificación. Una versión posterior del modelo podría
hacer depender la probabilidad de error de la longitud de la cola de triage,
lo que permitiría capturar un segundo mecanismo de mejora.

